\section{Memory-Augmented Manipulation Policies}
\label{sec:method}

Building on \benchname, we construct a family of memory-augmented vision-language-action (VLA) models based on the $\pi_{0.5}$ backbone, collectively termed the \textbf{\modelname\xspace suite}.
We systematically compare different memory representations and integration mechanisms under controlled settings, as illustrated in \cref{fig:memory_design}.
More detailed model formulations are provided in \cref{sec:appx_model_design}.

\subsection{Memory Representations}

We study three representations: \textit{symbolic}, \textit{perceptual}, and \textit{recurrent}.
Symbolic memory uses discrete language tokens, while the latter two provide differentiable neural features.

\noindent{\bf Symbolic Memory} is represented as interpretable language subgoals, as correct subgoals imply effective history summarization. At each step, we leverage an auxiliary vision-language model (VLM) that conditions on the current image and prior subgoals to generate the next subgoal. We consider \emph{simple} instructions and \emph{grounded} variants that additionally include image coordinates (e.g., \textit{pick up the green cube} vs. \textit{pick up the green cube at [63, 152]}), with grounding being particularly beneficial for spatial reasoning. Notably, symbolic memory approaches typically require additional subgoal annotations to train the VLM.

\noindent{\bf Perceptual Memory} is represented as a sequence of visual tokens selected from past images and  extracted by the $\pi_{0.5}$ vision encoder. To select informative tokens, we adopt two strategies: (1) token dropping \cite{yao2025timechat}, which removes temporally redundant image patches based on RGB differences; (2) uniform frame sampling \cite{chen2024videollm}, which evenly downsamples the sequence and concatenates tokens from the sampled frames. In practice, we found that visual tokens alone suffice to encode history without proprioceptive states.

\noindent{\bf Recurrent Memory} compresses the visual token sequence into fixed-size latent states through recurrence. We adopt two models: (1) Test-Time Training (TTT) \cite{sun2024learning, zhang2025test}, which updates fast weights online via a self-supervised loss and applies them to generate output features; (2) Recurrent Memory Transformers (RMT) \cite{bulatov2022recurrent, bulatov2023scaling}, which process input sequence into segments and recurrently update a set of learnable memory slots for each segment using transformer \cite{vaswani2017attention}.


\subsection{Memory Integration Mechanisms}
Symbolic memory can be naturally integrated into $\pi_{0.5}$  by concatenating subgoals with task instructions. In contrast, perceptual and recurrent memory produce neural \emph{memory tokens} that require more architectural modifications. We therefore study three integration mechanisms.


\noindent\textbf{Memory-as-Context.}
Memory tokens are concatenated with the original inputs and jointly processed by the VLM expert, directly influencing VLM features.

\noindent\textbf{Memory-as-Modulator.}
We adopt adaptive LayerNorm (AdaLN) \cite{peebles2023scalable} to condition the action expert on external memory. Before entering the feed-forward block in each layer, the action features cross-attend to memory tokens via multi-head attention to extract memory-aware representations. These representations are then projected into scale and shift parameters that modulate the normalized action features through AdaLN.

\noindent\textbf{Memory-as-Expert.}
We introduce an additional lightweight memory expert that processes memory tokens through a dedicated pathway. The three experts interact via a specified blockwise causal attention scheme: the action expert attends to both the VLM and memory experts, while the VLM and memory experts do not attend to each other. This separation limits interference, preserves the original VLM behavior, and allocates specialized capacity for memory processing.
